% ===========================================================================
% Section 5: Methodology (~2 pages)
% ===========================================================================

% --- 5.1 Bottleneck Taxonomy ---
\subsection{Bottleneck Taxonomy}
\label{sec:method:taxonomy}

% TODO: Define the bottleneck dimensions
% Multi-dimensional I/O health profile (not fixed classes):
% - Access Granularity (small vs. large I/O operations)
% - Metadata Intensity (excessive open/close/stat calls)
% - Parallelism Efficiency (load imbalance across ranks)
% - Access Pattern (random vs. sequential)
% - Interface Appropriateness (POSIX vs. MPI-IO for parallel workloads)
% - File Strategy (many small files vs. few large files)
% - Throughput Utilization (achieved vs. peak bandwidth)
% - Temporal Pattern (I/O bursts, synchronization barriers)
%
% Each dimension: continuous severity [0-1] + binary threshold flag
% Final taxonomy determined empirically from benchmark data

% --- 5.2 ML Classifier ---
\subsection{ML Classifier}
\label{sec:method:ml}

% TODO: Describe model training
% - Multi-label classification (BCE loss)
% - Models compared: XGBoost, LightGBM, Random Forest, MLP
% - Input: 157 normalized features
% - Output: probability per bottleneck dimension
% - Training on gold + silver labels with sample weighting
% - Hyperparameter tuning: Optuna with 5-fold temporal CV
% - Metric: micro-F1 (primary), macro-F1, per-class AP

% --- 5.3 SHAP Attribution ---
\subsection{SHAP Attribution}
\label{sec:method:shap}

% TODO: Describe SHAP usage
% - TreeSHAP for tree-based models
% - Per-sample feature importance ranking
% - Top-K features used for RAG query construction
% - Enables interpretable bottleneck explanations

% --- 5.4 LLM Recommendation Engine ---
\subsection{LLM Recommendation Engine}
\label{sec:method:llm}

% TODO: Describe LLM integration
% - RAG pipeline: SHAP top-K features -> vector search in KB
% - KB entries: {Darshan signature, code pattern, verified fix, speedup}
% - Prompt template: bottleneck labels + SHAP features + KB matches
% - LLM generates: (1) root cause explanation, (2) code-level fix,
%   (3) expected improvement
% - Grounding: every recommendation traceable to benchmark evidence
% - Evaluation: human expert rating + benchmark re-run verification
